\section{Results}
\label{sec:results}

% TODO: fill in results after agreeing on analysis plan
% The following placeholders reflect the analyses the pipeline supports.

\subsection{Corpus and coding overview}
% Summary table: per substance, n trips, total word count, n coders, n scenes individuated, n shared.

\subsection{Inter-rater consistency}
% - Scene-individuation agreement per coder pair (Table).
% - Conditional classification agreement at taxonomy depths 0, 1, 2.
% - Key finding: 81% of disagreement is attribute-level (classification), 19% individuation.
% - Binning view: confirms adjudication already captured span-overlap cases.
% - Rater-style asymmetries: per-coder tag-usage counts.

\subsection{Phenomenological profile of psilocybin}
% Consensus-filtered scene and attribute frequencies.
% - dominant alteration type
% - dominant object classes in visual hallucinations
% - auditory / tactile / nonsensory modality frequencies
% - emotion and insight profiles
% - ego dissolution frequencies

\subsection{Phenomenological profile of brugmansia}
% same structure

\subsection{Direct comparison}
% - items significantly more frequent in one substance than the other
% - qualitatively characteristic differences:
%   - psilocybin: geometric patterns, ego dissolution, metaphysical insight, extraordinary entities
%   - brugmansia: incrusted object hallucinations of humans/acquaintances,
%     delusional belief states, amnesia
% - scene-individuation rates per substance (are brugmansia trips more
%   hallucination-dense than psilocybin trips per word, etc.)
